% !TeX root = ../../../geos_mpm_manual.tex
\section{Native axisymmetric mechanics, contact and cohesive zones}
\label{sec:axisymmetric}
\index{axisymmetry}
Set \mpmcode{formulation=Axisymmetric} in PFW. The same timestep uses the
native 2D mesh with $x=r$ and $y$ axial. The geometric z coordinate and
circumferential velocity are zero. Constitutive tensors use the
$(r,y,\theta)$ basis: their zz entry is the hoop component. G and $G_0$
remain two full meridional edge generators, for a parallelogram or triangle.
No extra radius, hoop generator or vertex array is stored.

\subsection{Full-ring measures and cylindrical mapping}
PFW writes \mpmcode{spatialDimension=2} and \mpmcode{axisymmetric=1} on
\texttt{InternalMesh}. The partition retains this flag on restart; solver and
mesh must agree. The physical radial minimum is nonnegative, all particle
centers have positive radius, and radial periodicity is forbidden. Standard
exterior mesh buffer cells remain; they are not physical negative-radius
material. The solver applies radial symmetry at $r=0$.

Use full-ring physical quantities, without \texttt{planeStrainThickness}:
\[
 V_p=2\pi r_p A_p,\qquad m_p=\rho_p V_p,\qquad
 V_\Omega=\pi(r_{\max}^2-r_{\min}^2)(y_{\max}-y_{\min}).
\]
The particle center is the meridional area centroid. PFW evaluates every
emitted child's ring volume at its own radius, including cut triangles and
Cartesian-to-CPTI tessellation. Particle-sized fracture lengths use
$\sqrt{V_p/(2\pi r_p)}$; scalar user-specified lengths are unchanged.

The radius-weighted operators are
\[
 S_I=\frac{\int_A r N_I\,dA}{A r_p},\qquad
 \mathbf G_I=\frac{\int_A r\nabla N_I\,dA}{A r_p},\qquad
 T_I=\frac{\int_A N_I\,dA}{A r_p}.
\]
CPDI/CPTI integrate the material-domain interpolant of the grid basis. The
otherwise inactive third mapping-gradient slot stores $T_I$, not a z
spatial derivative. Point Q1 uses $T_I=N_I/r_p$. CPDI has four corners;
CPTI has three. The cylindrical force and hoop kinematics are
\begin{align*}
 f_{Ir}&=-V_p(\sigma_{rr}G_{Ir}+\sigma_{ry}G_{Iy}+\sigma_{\theta\theta}T_I),\\
 f_{Iy}&=-V_p(\sigma_{ry}G_{Ir}+\sigma_{yy}G_{Iy}),&
 L_{\theta\theta}&=\sum_I v_{Ir}T_I.
\end{align*}
FLIP, PIC, XPIC and FMPM include the hoop term. F, physical volume and stress
power use all three tensor diagonal components; only the 2D block advances G.
Material symmetries/orientations must be compatible with no swirl.

Temporary integration geometry is uniformly shrunk about its center if a
corner would pass below $r=0$. Stored material geometry, mass and physical
volume are unchanged. A center reaching $r\leq0$ is an error. Splitting uses
area fraction times $r_{\mathrm{child}}/r_{\mathrm{parent}}$ for extensive
state, with a weighted-mean correction for optional affine child velocities.
This conserves momentum, but resolving the velocity gradient can add resolved
kinetic energy, as in the existing splitter. Automatic resets preserve hoop
stretch and use $J/F_{\theta\theta}$ for planar geometry projection.

\subsection{Supported scope and deferred integrations}
The implementation supports ordinary contact between separate body groups,
Q1 point, CPDI and CPTI particles, fixed Outflow/Symmetry boundaries, and
compatible body-force, initial-velocity and material initialization.
It supports DFG contact-surface insertion. It explicitly rejects finite-thickness
polymer CZ, surface tension, reconstructed surfaces/nodal areas,
SPH/overlap corrections, crack-tip detection, profileHistory slice normalization,
B-splines, prescribed BC/F tables
and stress-controlled domain motion. PFW rejects surface traction generation;
direct GEOS particle tractions are interpreted as already integrated full-ring
forces. The uniform-confinement follow-up supports
\mpmcode{ConfiningPressure} with \mpmcode{uniformPressure=1} and no box/sphere
exclusion. Its background stress includes the hoop component; pressure is
retained after the event and across restart. Spatially restricted pressure
remains unsupported. Most geometry-changing/healing events are rejected. The prescribed-interface
\mpmcode{ReferenceCohesiveZones} event is supported with explicit surface data.

Built-in area-only Weibull objects are rejected: a statistical cell needs
$2\pi\int_A r\,dA$, before minimum-volume clamping. A custom wrapper can
implement \texttt{setAxisymmetric()}; explicit StrengthScale values are not
rescaled. Plane-strain thickness binding remains unchanged. Weibull and
localized hoop-driven cracking are deferred because circumferential
localization requires a treatment beyond the current symmetry assumption.

Silo/VTK output continues to show meridional quad/triangle glyphs at z=0.
Stress zz is hoop stress; volume, mass and reactions are full-ring values.
Reaction-history length\_z is zero; do not use Cartesian thickness-based
force-to-pressure conversions.
No automatic revolved output mesh is generated. ParticleColor is retained.

\subsection{Example, validation and remaining work}
\path{examples/axisymmetricCylinder} supplies a CPDI core with CPTI outer
annulus and radial contraction pulse. Its radius-1, length-1 cylinder has
volume $\pi$ and mass $1000\pi$. It is a smoke test, not a validated dynamics
benchmark. The runner generates inputs and optionally runs a selected GEOS
executable; its particle check verifies geometry and full-ring volume.

Standalone checks cover 1800 production axisymmetric mappings, cylindrical
moments, hoop work, hydrostatic radial equilibrium, near-axis restriction and
split conservation. Complete writer checks cover both ordinary and curved
refined geometry at 1, 3 and 7 explicitly emulated ranks. These checks do not
replace the target GEOS build, actual MPI migration/restart, GPU execution or
output readback. See \path{AXISYMMETRIC_PLAN.md} for exact limits, local results,
and the remaining polymer CZ, traction and verification stages.

\path{examples/axisymmetricJacketedSample} adds a 20-mm-diameter, 20-mm-long
compliant sample between 2-mm steel plates inside a 1-mm rubber jacket.
CPDI quads fill the sample and plates; CPTI triangles fill the jacket and end
caps. A uniform pressure bath ramps to 10 kPa over 1 ms and holds for 1 ms.
All interfaces remain bonded in group zero. The sample is an effective elastic
porous continuum, without explicit pores, pore fluid, failure or plastic
compaction. An equal-material variant checks the small-strain homogeneous
reference $\sigma_{rr}=\sigma_{yy}=\sigma_{\theta\theta}=-p$ and
$\Delta V/V=-p/K$. Source-level checks do not establish dynamics convergence;
run the included histories and resolution checks on the target build.


\subsection{Ordinary contact between bodies}
For ordinary contact without DFG, assign separate nonnegative
\mpmcode{ContactGroup} values and set
\mpmcode{enableContact=1}, with \mpmcode{damageFieldPartitioning=0} and
\mpmcode{overlapCorrection=Off}. Existing normal-selection, gap-gating and
Coulomb-friction options are retained. Nodal mass and momentum already use
full-ring measures, so the contact force/impulse law needs no extra $2\pi r$.
The shared boundary-normal reconstruction is
\[
 \boldsymbol b_I=\sum_p V_p\bigl(G_{Ir}+T_I,\ G_{Iy},\ 0\bigr).
\]
The added radial term follows from
$\nabla(rN_I)=r\nabla N_I+N_I\boldsymbol e_r$ and removes the radial bias
that would otherwise remain inside a fully occupied nodal support. It does
not prescribe extra circumference weighting for supplied unit normals.
The same assembled normals serve ordinary and FMPM contact.

Explicit PFW surface positions can be used with
\mpmcode{useSurfacePositionForContact=1}; otherwise the existing approximate
centre-of-mass gap criterion applies. Solver-reconstructed surface positions,
overlap correction, nodal areas and moving walls remain
unsupported. Contact and slip have only radial and axial components.

\path{examples/axisymmetricContact} supplies two elastic cylinders with separate
contact groups and mixed CPDI/CPTI particles. Each has radius 10 mm and length
8 mm; they approach at $\pm1$ m/s across a 2-mm gap. The example includes input
checks, a runner and history plots. \mpmcode{AXIS_CONTACT_UPDATE=FMPM} selects
FMPM(2); the default is FLIP. Standalone tests check normal cancellation and
orientation on axial, cylindrical and inclined interfaces, axis behavior,
friction limits, separation, momentum exchange and contact dissipation. Writer
checks use emulated collectives; full GEOS contact dynamics and MPI/restart
validation remain target-machine tasks.

The uniform pressure bath used in the bonded jacketed example can also load
newly exposed internal surfaces if the materials are separated into contact
groups. A sealed jacket with internal separation needs exterior-only pressure
loading, which remains a separate integration task.


\subsection{Axisymmetric DFG contact-surface insertion}
Set \mpmcode{damageFieldPartitioning=1} and \mpmcode{enableContact=1}.
The existing DFG field assignment, separability law and defaults remain in
use. For $n_g$ contact groups, DFG supplies two slots per group with field
index $\mathrm{side}\,n_g+\mathrm{group}$. Ordinary contact evaluates the
occupied field pairs in both the standard and FMPM paths. This does not
create cohesive bonds or generalize the CZ field-pair representation.

Damage is a scalar on the meridional section. Its reconstruction uses
\[
 A_q=\frac{V_q}{2\pi r_q},\qquad
 \bar d(\boldsymbol x)=\frac{\sum_q A_q d_q K(\boldsymbol x-\boldsymbol x_q)}
                              {\sum_q A_q K(\boldsymbol x-\boldsymbol x_q)},\qquad
 \nabla\bar d=\frac{\sum_q A_q(d_q-\bar d)\nabla K_q}{\sum_q A_q K_q}.
\]
Use the existing compact cubic kernel and neighbor set. Each sample also
contributes its image at $(-r_q,y_q)$ when within support. This even radial
extension makes the radial derivative zero at the axis; it creates no
particles or additional physical mass. The existing surface-as-damage
convention is retained. Particle positions and volumes are current values.
Only this reconstruction changes measure: mapped damage, quality tensors,
volume fractions, nodal mass, momentum and contact forces retain full-ring
weights. No additional circumference factor enters the contact law.

The existing separability criterion decides whether same-group fields remain
tied or can open and recontact. Its development is separate from this
axisymmetric adaptation. Solver-reconstructed surface positions remain
unsupported; a newly inserted DFG surface can use the existing approximate
centre-of-mass gap criterion with \mpmcode{useSurfacePositionForContact=0}.
The current partition uses one reference direction per grid node and two
slots per group; this patch does not add unresolved crack-junction topology.
Only cracks whose revolution is axisymmetric are representable. Localized
longitudinal hoop fractures and swirl remain outside the formulation.

\path{examples/axisymmetricDFG} contains four isolation cases: an intact
cylinder, the same cylinder with a prescribed damaged central band, three
approaching contact groups, and the three groups with damage in the central
body. The first pair initially moves away from the midplane. There is no
applied tensile end traction, new damage-evolution law, explicit internal
normal, or CZ activation. Select with \mpmcode{AXIS_DFG_CASE}; use
\mpmcode{AXIS_DFG_UPDATE=FMPM} for FMPM(2), otherwise FLIP. Particle checks,
runner and tracer postprocessing are included. See \path{AXISYMMETRIC_DFG.md}
for the standalone kernel and writer checks and the remaining target tests.

\subsection{Prescribed axisymmetric cohesive interfaces}
Use exactly two contact groups, 0 and 1, and disable DFG. A
\mpmcode{CohesiveZoneRegion} selects particles through its tag; surface
particles use \mpmcode{SurfaceFlag=3}. PFW must include ContactGroup, CZTag,
SurfaceFlag, SurfaceNormal and SurfacePosition columns. The surface vectors
are meridional and the normal is nonzero. Set \mpmcode{useEvents=1} on the
solver and activate the region with
\mpmcode{ReferenceCohesiveZones}, using \mpmcode{computeNormalsAndPositions=0}.
The event dispatcher defaults to disabled; declaring the event and marking
surface candidates alone does not set active particle cohesive flags.
Each reference event requests initialization once, using an internal
\mpmcode{referenceApplied} flag saved across restart. A time window alone is
insufficient: timestep growth can admit repeated activation and overwrite the
bonded reference after separation begins. Dependency scheduling is distinct
from activation, so its \mpmcode{hasStarted} flag is not used for this purpose.
Finite event end times still disable the region; a separate reference event
can intentionally initialize it again. This fix leaves general event
scheduling unchanged and applies to all formulations. Rebuild and restart a
case affected by repeated initialization from time zero.
Uncoupled, Coupled and Bicrystal traction--separation laws are supported.
Finite-thickness Polymer and SurfaceInformedPolymer CZ models remain rejected:
they need hoop stretch in their constitutive deformation reconstruction.

At activation, the reference nodal area is
\[
 A_{I,\mathrm{CZ}}=\int_{\Gamma}2\pi r N_I\,ds.
\]
The reconstructed local line is clipped to nodal support and $r\geq0$, split
at Q1 hat-function kinks, and integrated exactly with two-point Gauss rules.
The optional region \mpmcode{czVolumeNormalization=1} uses the matching support
volume $\int 2\pi r N_I\,dA$; at an axis node it is
$2\pi h_r^2h_y/6$. Multiplication by the node radius alone is incorrect.
Finite interface ends retain the existing local-line occupancy approximation;
verify area and reaction convergence with refinement.

Current area uses the full three-dimensional relative cofactor,
\[
 H=F F_{\mathrm{CZ}}^{-1},\qquad
 \mathbf a_I=\langle\operatorname{cof}H\rangle_I
              A_{I,\mathrm{CZ}}\mathbf n_{I,\mathrm{CZ}}.
\]
Thus axial and cylindrical areas include hoop stretch. Reference normals are
normalized. Re-basing $F_{\mathrm{CZ}}$ during an F reset preserves H and area.
Both FOnly and FAndGeometry retain the cached cohesive anchor. TypeB selects
that anchor by intersecting the normal ray with the unscaled quad/triangle
material domain; TypeA uses the supplied surface offset and Nodal retains the
reference-grid affine reconstruction. These choices change the displacement
reconstruction, not the reference-area integration.

The law uses opening and meridional slip; force is traction times the current
full-ring area once. With \mpmcode{preventCZInterpenetration=1}, ordinary contact
handles compression. A particle retains cohesive forces while any cached CZ
node remains intact; complete failure releases it to ordinary contact.
Contact handoff remains particle-based, not an independently tracked per-node
crack front. Active/damaged cohesive surface particles are not subdivided.
Scaling/reset options preserve their surface normals and positions.

\mpmcode{gridCohesiveArea} now contains current area vectors, and
\mpmcode{gridCohesiveForce} contains forces, at cached reference-grid locations.
Meridional Silo/VTK output needs no extra thickness conversion.
\path{examples/axisymmetricCZ} provides mixed CPDI/CPTI elastic and debond
variants, a runner, input checks and tracer plots. Standalone tests compile
production area, CZ-law, relative-F, reset and gather excerpts. An event test
also compiles the reference event and scheduler/handler excerpts to check
adaptive timesteps, dependency delays and restart-wrapper state. Writer tests
exercise FLIP/FMPM with one and three emulated ranks. Full GEOS dynamics,
actual MPI/restart and GPU execution remain target tests. See
\path{AXISYMMETRIC_CZ.md} for the precise validation scope.
